% !PRINTED 390
\weberpage{390}
%<<CONT
a number in $A_1$ with a number in $A_2$, which we express by

\weq{2}{Z = A_1 A_2}

in symbolic form.

The \gesperrt{intersection} $A$ of the groups $A_1$ and $A_2$ contains all and only those numbers which are at the same time norm residues of $m_1$ and of $m_2$, and we prove that the group of the norm residues $m$

\weq{3}{m = m_1 m_2}

is $A$.

For if $a$ is a number which is contained at the same time in $A_1$ and in $A_2$, then there are two numbers $\omega_1, \omega_2$ in $O$ for which the conditions

\weq{4}{\begin{aligned} N(\omega_1) &\equiv a \pmod{m_1},\\ N(\omega_2) &\equiv a \pmod{m_2} \end{aligned}}

are satisfied. Two rational numbers $x_1, x_2$ can now be determined in such a way that

\[
\begin{aligned}
x_1 &\equiv 1 \pmod{m_1}, & x_2 &\equiv 0 \pmod{m_1},\\
&\equiv 0 \pmod{m_2}, & &\equiv 1 \pmod{m_2},
\end{aligned}
\]

and if we then accordingly put

\[ \omega = x_1 \omega_1 + x_2 \omega_2 \]

we have

\weq{5}{\begin{aligned} N(\omega) &\equiv N(\omega_1) \equiv a \pmod{m_1},\\ &\equiv N(\omega_2) \equiv a \pmod{m_2}, \end{aligned}}

and since $m_1$ and $m_2$ are relatively prime, we also have

\weq{6}{N(\omega) \equiv a \pmod{m},}

so that $a$ is a norm residue of $m$. Conversely, a norm residue of $m$ is at the same time a norm residue of every divisor of $m$, and hence also of $m_1$ and $m_2$.

From this there now follows, by means of Theorems 13 and 14, \S\,100:

\[
\begin{aligned}
(Z, A_1) &= (A_1 A_2, A_1) = (A_1 A_2, A_1 A) = (A_2, A),\\
(Z, A) &= (Z, A_2)(A_2, A) = (Z, A_1)(Z, A_2),
\end{aligned}
\]

and we have the relation:

\weq{7}{(Z, A) = (Z, A_1)(Z, A_2).}

With this the question of the norm residues is reduced to the case in which the modulus is a \gesperrt{prime power}.

\S\,108. \gesperrt{Norm residues of prime powers.}

Let $O$ be an order group in the field $\Omega$ (\S\,98) and $P$, the groups of the rational numbers which are relatively prime to the discri-
%CONT>>
